% =============================================================================
% 04C_REGULATORY_COMPETITIVE_ANALYSIS — Comparative Analysis of Regulatory
% Frameworks Against the Three-Gap Framework
% =============================================================================
% Compares EU AI Act, US FDA, UK MHRA, and WHO frameworks on their capacity
% to address the gap1 → gap2 → gap3 cascade identified in the empirical analysis.
% =============================================================================

\section{Regulatory Competitive Analysis}
\label{sec:regulatory}
\addcontentsline{toc}{section}{Regulatory Competitive Analysis}

The empirical analysis in Section~\ref{sec:empirical} identified the
\textbf{gap1 $\rightarrow$ gap3 cascade} as the dominant causal pathway
in the literature: ecological validity failures in evaluation drive pilot-trap
outcomes. Gap 1 is the upstream regulatory lever with the largest downstream
effect on implementation success. This section applies that finding to a
comparative analysis of the four major regulatory frameworks operating in
the health AI space: the \textbf{EU AI Act}, the \textbf{US FDA framework},
the \textbf{UK MHRA framework}, and the \textbf{WHO Governance Guidance}.
Each framework is evaluated against its capacity to address the three-gap
cascade, with particular attention to the gap1 $\rightarrow$ gap3 pathway
that the empirical record identifies as primary.

\medskip
\noindent\textbf{Methodological note.} The evaluative criteria applied in this
section (the five conditions derived from the three-gap framework) are
empirically grounded in the 60-paper v3 corpus. However, the application of
those criteria to specific regulatory frameworks draws on regulatory knowledge
that extends beyond the strict corpus. This is methodologically appropriate for
a literature review: the corpus identifies \textit{what the gaps are} and
\textit{how they interact}, while the regulatory comparison applies those
findings to the broader governance landscape. Where specific regulatory claims
are made, they are supported by citations to corpus papers where available
(e.g., \cite{warraich2024fda, abulibdeh2025illusion} for FDA critique) and
to authoritative regulatory documentation where the corpus does not provide
direct evidence. The reader should therefore understand this section as an
analytical application of the empirical framework to the known regulatory
landscape, not as a systematic content analysis of the 60-paper corpus
against each regulatory framework.

\subsection{Evaluative Framework: What Should a Health AI Regulatory Framework Address?}

Before comparing frameworks, it is necessary to establish the evaluative
criteria derived from the three-gap framework. A regulatory framework that
genuinely addresses the health AI evaluation-governance-implementation
challenge must satisfy five conditions:

\begin{enumerate}
  \item \textbf{Ecological validity mandate:} The framework must require
        evaluation methodologies that reflect real-world clinical complexity —
        including external validation, diverse population subgroups, and
        longitudinal performance monitoring — not merely technical accuracy
        on curated single-site datasets.

  \item \textbf{Governance capacity recognition:} The framework must recognize
        that governance capacity is built through implementation, not pre-installed
        prior to deployment. Regulatory design must account for organizational
        readiness as a dynamic process rather than a static precondition.

  \item \textbf{Pilot-to-routine transition support:} The framework must
        provide mechanisms that actively support the transition from pilot
        to routine care — including staged authorization, enhanced monitoring
        during the transition period, and clear escalation pathways.

  \item \textbf{Algorithmic drift and post-market monitoring:} The framework
        must mandate continuous post-market performance monitoring with
        standardized drift detection protocols, moving beyond one-time
        pre-market authorization.

  \item \textbf{Equity and subgroup performance:} The framework must require
        mandatory subgroup performance reporting to detect and quantify
        algorithmic bias across population strata, addressing the equity gap
        documented in the literature.
\end{enumerate}

These five conditions operationalize the three-gap framework as a regulatory
design criterion. Each framework is assessed against all five.

% =============================================================================
% EU AI ACT
% =============================================================================

\subsection{EU AI Act}
\label{sec:reg_eu}

The European Union's AI Act \citep{eu_ai_act_2024}, which entered into force
in August 2024, represents the world's most comprehensive horizontal
regulation of AI systems. Its risk-based architecture classifies AI systems
into four tiers — unacceptable risk (prohibited), high risk, limited risk,
and minimal risk — with health AI systems predominantly falling into the
high-risk category.

\textbf{Strengths against the three-gap framework.}

The EU AI Act scores highest on Condition2 (governance capacity recognition).
Article 10 requires that high-risk AI systems be accompanied by technical
documentation demonstrating not only technical performance but also
organizational readiness to manage the system throughout its lifecycle.
Crucially, Article 14 on human oversight mandates that high-risk AI systems
be designed to allow human intervention and that organizations maintain
processes for ongoing monitoring — implicitly acknowledging that governance
capacity must be maintained, not merely asserted at authorization.

The Act also scores well on Condition 4 (post-market monitoring). Article 10
requires manufacturers to implement quality management systems that include
post-market surveillance procedures, and Article 92 mandates that the European
AI Office conduct ongoing conformity assessments. This represents a meaningful
departure from the one-time authorization model.

On Condition 5 (equity and subgroup performance), the Act's requirements for
training data documentation (Article 10) and the risk management system
(Article 9) create institutional hooks for equity auditing, though mandatory
subgroup performance reporting in deployment is not explicitly required.

\textbf{Weaknesses against the three-gap framework.}

The EU AI Act scores poorly on Condition 1 (ecological validity mandate).
The conformity assessment procedure for high-risk AI systems (Annexes VII–VIII)
focuses on technical documentation and testing protocols but does not mandate
external, multicenter, prospective evaluation as a condition of market entry.
The Act's evaluation requirements are process-oriented — demonstrating that
a quality management system exists — rather than outcome-oriented in terms
of ecological validity. A manufacturer could satisfy the conformity assessment
with single-site retrospective validation data, which the empirical analysis
identifies as the primary source of the ecological validity gap.

On Condition 3 (pilot-to-routine transition), the Act provides no staged
authorization mechanism. High-risk AI systems are subject to the same
requirements at the pilot stage as at full deployment, which may create
disincentives for organizations to run controlled pilot phases with enhanced
monitoring. The regulatory sandbox provisions (Article 53) are designed for
innovation support but are not explicitly tied to a structured pilot-to-routine
transition pathway.

\textbf{Summary assessment for EU AI Act.} The Act is the most structurally
sophisticated of the four frameworks, with meaningful post-market monitoring
requirements and explicit governance capacity provisions. Its primary weakness
is the absence of mandatory ecological validity standards for pre-market
evaluation — it does not close the gap1 $\rightarrow$ gap3 cascade at its
upstream source. Organizations seeking to deploy AI under the EU AI Act
will encounter robust governance requirements but no regulatory pressure to
adopt ecologically valid evaluation methodologies.

% =============================================================================
% US FDA
% =============================================================================

\subsection{US FDA Framework}
\label{sec:reg_fda}

The US Food and Drug Administration's approach to AI-enabled medical devices
operates primarily through the 510(k) and De Novo clearance pathways, with
Software as a Medical Device (SaMD) governed under the FD\&C Act framework
\citep{warraich2024fda}. As of 2024, the FDA had authorized nearly 1,000
AI-enabled medical devices, making it the largest single regulatory authority
for clinical AI in the world.

\textbf{Strengths against the three-gap framework.}

The FDA scores highest on Condition 3 (pilot-to-routine transition support)
through its Predetermined Change Control Plan (PCCP) mechanism. Introduced
in 2023 and formalized in guidance, the PCCP allows manufacturers to define
in advance the types of model modifications that can be implemented without
requiring new market authorization, subject to documented performance
monitoring. This is a genuinely innovative mechanism: it creates a regulatory
structure for iterative deployment that mirrors the pilot-to-routine transition
process, allowing organizations to move from controlled updates to full
routine operation within a single authorization.

The FDA also scores reasonably well on Condition 4 (post-market monitoring).
The agency has issued guidance on real-world performance monitoring for
AI/ML-enabled devices and requires adverse event reporting through the
MAUDE database. The 2024 guidance on predetermined change control plans
further incorporates mandatory performance monitoring thresholds.

\textbf{Weaknesses against the three-gap framework.}

The FDA scores most poorly on Condition 1 (ecological validity mandate).
\citet{abulibdeh2025illusion} document that the majority of FDA-authorized
AI devices were approved on the basis of single-site, premarket studies
with no mandatory external validation. The 510(k) pathway in particular
operates on a substantial equivalence model that does not require prospective,
multicenter, ecologically grounded evaluation. The FDA's focus is on
demonstrating that a device is substantially equivalent to a predicate
device — a process-oriented standard — rather than on real-world performance.

On Condition 2 (governance capacity recognition), the FDA framework assumes
that organizations seeking authorization possess the governance infrastructure
necessary for safe deployment. The FDA does not assess organizational
governance maturity as a condition of authorization, and the guidance on
governance capacity building is minimal. This is the assumption that
\citet{tian2026pilottrap} identify as empirically challengeable.

On Condition 5 (equity and subgroup performance), the FDA has historically
not required mandatory subgroup performance reporting in authorization
submissions. While the 2023 action plan on AI/ML in medical devices
identifies equity as a priority, enforceable subgroup reporting requirements
remain under development.

\textbf{Summary assessment for FDA.} The FDA is the world's most active
regulator of clinical AI by volume, but its framework is the least well
adapted to the ecological validity challenge identified in the empirical
analysis. The PCCP mechanism is a genuine innovation for pilot-to-routine
transition, but the absence of mandatory ecological validity standards
means that the gap1 $\rightarrow$ gap3 cascade is not addressed at its
upstream source. The FDA framework generates authorization for AI systems
whose validation evidence systematically understates real-world complexity.

% =============================================================================
% UK MHRA
% =============================================================================

\subsection{UK MHRA Framework}
\label{sec:reg_mhra}

The UK Medicines and Healthcare products Regulatory Agency (MHRA) has developed
a distinctive approach to AI-enabled medical devices following Brexit,
positioning the UK as a potential regulatory innovator through its Software
and AI Incident Management Programme and the emerging AI roadmap
\citep{mhra2024software, mhra2024ai}. The MHRA's approach is notable for
its emphasis on iterative, real-world evidence and its use of the Yellow Card
reporting system for post-market surveillance.

\textbf{Strengths against the three-gap framework.}

The MHRA scores highest on Condition 4 (post-market monitoring) through its
established Yellow Card pharmacovigilance-style reporting system, which has
been extended to AI-enabled medical devices. The agency has also signaled
intent to develop real-world performance monitoring frameworks that go beyond
adverse event reporting to include proactive drift detection.

On Condition 3 (pilot-to-routine transition), the MHRA's emphasis on iterative
authorizations and its work on Software as a Medical Device (SaMD) pathways
suggests a framework better adapted to the iterative deployment lifecycle
than the FDA's 510(k) pathway. The MHRA's 2024 software and AI work programme
explicitly addresses the challenge of continuous learning AI systems.

\textbf{Weaknesses against the three-gap framework.}

The MHRA scores poorly on Condition 1 (ecological validity mandate). The
current MHRA guidance on AI device evaluation does not mandate external,
multicenter, prospective evaluation as a condition of authorization. The
framework relies on existing conformity assessment pathways that do not
require ecologically grounded evaluation.

On Condition 2 (governance capacity recognition), the MHRA framework
similarly assumes organizational readiness without assessing it as a
dynamic process. No guidance currently addresses governance capacity building
as a regulatory requirement.

On Condition 5 (equity and subgroup performance), the MHRA has not issued
enforceable subgroup performance reporting requirements, though this is an
area of active policy development.

\textbf{Summary assessment for MHRA.} The MHRA framework is the least
documented of the four in the reviewed literature, making definitive assessment
difficult. Its real-world evidence orientation and iterative authorization
approach suggest potential to address the pilot-to-routine transition, but
the absence of mandatory ecological validity standards and governance capacity
provisions limits its effectiveness against the three-gap cascade.

% =============================================================================
% WHO GOVERNANCE GUIDANCE
% =============================================================================

\subsection{WHO Governance Guidance}
\label{sec:reg_who}

The World Health Organization's guidance on ethics and governance of AI in
health \citep{who2024governance} represents the most authoritative global
normative framework for health AI governance. Unlike the EU AI Act, FDA, or
MHRA frameworks — which are legally binding within their jurisdictions —
the WHO guidance functions as soft law: it establishes normative standards
that member states are encouraged to adopt, but it does not carry enforcement
mechanisms.

\textbf{Strengths against the three-gap framework.}

The WHO guidance scores highest on Condition 2 (governance capacity
recognition) and Condition 5 (equity and subgroup performance). The guidance
explicitly addresses governance capacity as a dynamic, building process:
it advocates for continuous compliance mechanisms that evolve with the AI
system's performance history, and it places equity and inclusion at the
center of governance design. The guidance's emphasis on participatory
governance — incorporating affected communities, patient advocacy groups,
and frontline clinical staff — directly addresses the equity dimension
that the empirical analysis identifies as a gap.

On Condition 3 (pilot-to-routine transition), the WHO guidance recommends
phased deployment approaches with mandatory performance thresholds at each
stage, which is conceptually aligned with the pilot-to-routine transition
challenge.

\textbf{Weaknesses against the three-gap framework.}

The WHO guidance scores poorly on Condition 1 (ecological validity mandate).
As a normative document, the guidance does not specify evaluation methodology
requirements and does not mandate external validation or subgroup performance
reporting as technical standards. The guidance addresses governance and ethics
but not the technical evaluation methodology that the empirical analysis
identifies as the upstream driver of the cascade.

On Condition 4 (post-market monitoring), the WHO guidance advocates for
continuous monitoring but provides no operational mechanism, enforcement
pathway, or threshold standard. The advocacy is normative rather than
operational.

\textbf{Summary assessment for WHO.} The WHO guidance is the strongest
of the four frameworks on governance capacity recognition and equity, but
its lack of enforcement mechanisms and absence of ecological validity
mandates limit its effectiveness as a regulatory instrument. It functions
best as a complement to legally binding frameworks (EU AI Act, FDA, MHRA)
rather than as a standalone regulatory mechanism.

% =============================================================================
% COMPARATIVE SUMMARY
% =============================================================================

\subsection{Comparative Summary}
\label{sec:reg_comparison}

Table~\ref{tab:reg_comparison} summarizes the four frameworks against the
five evaluative conditions derived from the three-gap framework.

\begin{table}[t]
  \centering
  \caption{Comparative assessment of four regulatory frameworks against the
  five conditions derived from the three-gap framework. Ratings: ++ strong,
  + moderate, $-$ weak, $--$ absent or counter-productive.}
  \label{tab:reg_comparison}
  \begin{tabular}{lcccc}
    \toprule
    \textbf{Condition} & \textbf{EU AI Act} & \textbf{FDA} & \textbf{MHRA} & \textbf{WHO} \\
    \midrule
    Ecological validity mandate (G1)      & $-$  & $--$ & $-$  & $--$ \\
    Governance capacity recognition (G2)  & ++   & $-$  & $-$  & ++ \\
    Pilot-to-routine transition (G3)       & $-$  & ++   & +    & + \\
    Post-market monitoring (drift)         & ++   & +    & +    & $-$ \\
    Equity and subgroup reporting          & +    & $-$  & $-$  & ++ \\
    \midrule
    \textbf{Addresses gap1$\rightarrow$gap3 cascade}
                                           & Partial & Partial & Partial & No \\
    \bottomrule
  \end{tabular}
\end{table}

The central finding of this comparison is that \textbf{no existing framework
adequately addresses the gap1 $\rightarrow$ gap3 cascade}. All four frameworks
share a common structural weakness: they address governance capacity (Gap 2)
and implementation barriers (Gap 3) with varying degrees of effectiveness,
but none mandates ecologically valid evaluation methodologies as a condition
of market authorization. This is the upstream gap that the empirical analysis
identifies as the primary driver of the cascade.

Three patterns emerge from this comparison:

\begin{enumerate}
  \item \textbf{The ecological validity gap is regulatory blind spot.}
        Across all four frameworks, the absence of mandatory external,
        multicenter, prospective evaluation requirements means that
        AI systems enter the market with validation evidence that
        systematically understates real-world complexity. Regulatory
        frameworks are designed to ensure that marketed products meet
        defined standards — but those standards do not currently include
        ecological validity. Closing this gap requires adding ecological
        validity requirements to the conformity assessment procedures of
        all four frameworks.

  \item \textbf{Governance capacity is increasingly recognized but not
        operationalized.} The EU AI Act and WHO guidance both explicitly
        address governance capacity as a dynamic process, but neither
        provides operational tools for assessing organizational readiness
        in real time. Hussein's maturity model \citep{hussein2026governance}
        and Alsaigh's PEARL-PATHWAY framework \citep{alsaigh2026pearl}
        offer the operational instruments that regulatory frameworks
        currently lack. The gap between regulatory recognition and
        operational implementation of governance capacity is where the
        pilot trap takes root.

  \item \textbf{The PCCP mechanism is the most promising regulatory
        innovation for addressing Gap 3.} The FDA's Predetermined Change
        Control Plan is the only mechanism among the four frameworks that
        creates a regulatory structure explicitly designed for the
        pilot-to-routine transition. Its staged authorization model —
        with mandatory performance thresholds at each stage — is the
        closest approximation to the kind of adaptive regulatory approach
        that the empirical analysis calls for. The EU AI Act's sandbox
        provisions are conceptually similar but less operationally
        specified.
\end{enumerate}

The implication for regulatory design is clear: closing the gap1
$\rightarrow$ gap3 cascade requires adding an ecological validity mandate
to existing regulatory frameworks, operationalizing governance capacity
assessment through maturity models and readiness tools, and extending
the PCCP-style staged authorization model across all jurisdictions.
The three-gap framework provides the analytical structure for identifying
where regulatory intervention is most needed — and the empirical analysis
identifies ecological validity as the upstream lever with the largest
downstream effect on the cascade.
